% 05_upper.tex
\section{The upper side and the two-sided theorem}
\label{sec:upper}

\begin{lemma}[closed-form upper bounds of the branches] \proofst
\label{lem:AZ1P}
On every periodic window,
$(a)\le2\sigma^2J_{ww,\mathrm{up}}/(\sigma\,\mathrm{den}_{\mathrm{lo}})$ with $J_{ww,\mathrm{up}}=(1+u_0/(8(1-\Qb)))/L^2$ and
$\mathrm{den}_{\mathrm{lo}}=b+xI_2/(2(1+\Qb))$; and $(a')\le\tfrac{\sigma^2}{2}\E[s_v^2\phi]/(\sigma\,\mathrm{den}_{\mathrm{lo}})$ with
\begin{equation}
\begin{aligned}\E[s_v^2\phi]\le{}&(1-e)J_{\mathrm{per,up}}+\tfrac14\langle\tilde s^2R^2\rangle\\&+\frac{1-e}{V}+\frac{1-C}{V}+\alpha^2eC+\gamma_P(1-C),\end{aligned}
\end{equation}
$1-C\le(L+r)^2/4$, $J_{\mathrm{per,up}}=xI_2/(2(1-\Qb)\sigma^2)$, $e=e^{-\alpha^2V/2}$, $\alpha=(1-\sigma^2/V)/\sqrt2$, and
$\langle\tilde s^2R^2\rangle$, $r$ bounded by Minkowski and the Gaussian moments of the common envelope (App.~\ref{app:plemmas}).
With~\eqref{eq:branch}: $g(\rho)\le(1+|\varepsilon(\rho)|)U_P(\rho)$, $U_P:=(a)_{\mathrm{up}}+(a')_{\mathrm{up}}$, for every periodic window $\rho$.
(The averages here are period averages; Lemma~\ref{lem:BL7} gives the $\theta_\tau$ remainders, which along the comb sequence below are smaller
than $e^{-10^{22}}$ relative to $U_P$, since there $\tau=k'^2V\ge8\times10^{22}$.)
\end{lemma}

\paragraph*{Two explicit sequences.}
The strip sequence $S'$: $c=1$, $4k_0^2+11\ln k_0=\ls$, $M=k_0^2e^{k_0^2}$, $\eta=1/k_0^2$, $\lambda=2/\sqrt{k_0}$. The comb sequence $S'_C$ is its image
under the map of Proposition~\ref{prop:BE1} with $a=2$:
\begin{equation}
\begin{aligned}&4k_0^2+11\ln k_0=\ln(256/\sigma),\\ &M=\tfrac14k_0^2e^{k_0^2},\quad \eta=\tfrac1{k_0^2},\quad\lambda=\tfrac2{\sqrt{k_0}},\end{aligned}
\end{equation}
with Gaussian teeth $a_n=2e^{-n^2/(2k_0^4)}$ ($a_1\to2$, $|a_n|\le a_1$, $f$ even). Along both sequences every length is an exact monomial in
$k_0$ and $e^{k_0^2}$ (for $S'_C$: $\sqrt\sigma=16k_0^{-11/2}e^{-2k_0^2}$, $\ell=32k_0^{-6}e^{-2k_0^2}$, $L_w=8k_0^{-4}e^{-k_0^2}$,
$L_v=32k_0^{-8}e^{-2k_0^2}$), which is what makes a rigorous evaluation possible far below double precision.

\begin{theorem}[upper side on the comb sequence] \caproofst{\code{bf\_cert\_upper.py}, \code{bs\_upper\_tail.py}}
\label{thm:AZ1P}
Along $S'_C$, $g(S'_C(\sigma))\le(1+|\varepsilon(S'_C(\sigma))|)\,U_P(\sigma)$ and
\begin{equation}
U_P\sqrt{\ls}\ \le\ 1+6.0\,\frac{\ln\ls}{\ls}\qquad(0<\sigma\le10^{-24}).
\end{equation}
Moreover $U_P\sqrt\ls\to1$ and $\limsup_{\sigma\to0}(U_P\sqrt\ls-1)\ls/\ln\ls\le\frac{11}4$.
The same holds along the strip sequence $S'$ with the constant $7.4$ (certified in the same way, \code{ba\_cert.py}).
\end{theorem}

The one-variable inequality is verified on $1464$ cells covering $k_0\in[3.43,240]$ ($\ls$ from $55.07$ to $230455$) with exact rational
exponents for the large factors and $60$-digit outward-rounded interval arithmetic for the $O(1)$ factors; the harmonics $n\ge2$ enter
through interval parameters of width $\le6\times10^{-7}$ (Appendix~\ref{app:cert}). For $k_0\ge240$ (v1.3) every factor of $U_P$ is written as its
leading monomial in $k_0$ and $e^{-k_0^2}$, with an exact leading constant, times $(1+R\,k_0^{-\nu})$ with an explicit interval $R$; this gives
$U_P=\frac1{2k_0}(1+R_U/k_0^2)$ with $R_U\le3.5356$ for all $k_0\ge240$ at once, and the inequality follows from an elementary bound whose margin increases
with $k_0$ (App.~\ref{app:cert}). Since $\ls=4k_0^2+O(\ln k_0)$, the same bound gives $U_P\sqrt\ls\to1$. The comb's constant is smaller than the strip's because
$S'_C$ is $S'$ at $\sigma/256$ and $\ls$ is shorter by $\ln256$. Paper~E's evaluator confirms the bound (Table~\ref{tab:upper}): it holds
for $(a)$, $(a')$ and $g$ separately and becomes sharp as $\sigma\to0$ \numst.

\begin{table}
\caption{The comb sequence $S'_C$ against paper~E's evaluator (four harmonics; bound with the same four).}
\label{tab:upper}
\begin{ruledtabular}
\begin{tabular}{lccccc}
$\sigma$ & $a_1$ & $\Qb$ & $g$ & $U_P$ & $U_P/g-1$\\ \hline
$10^{-6}$ & 1.906 & 0.38 & 0.6323 & 1.1452 & 0.81\\
$10^{-12}$ & 1.971 & 0.11 & 0.3405 & 0.4042 & 0.19\\
$10^{-18}$ & 1.987 & 0.025 & 0.2401 & 0.2506 & 0.043\\
$10^{-24}$ & 1.993 & 0.0054 & 0.1904 & 0.1923 & 0.0097\\
\end{tabular}
\end{ruledtabular}
\end{table}

\subsection{The two-sided theorem}

\begin{table}
\caption{The assumptions of Theorem~\ref{thm:BF1}, and what they exclude.}
\label{tab:assumptions}
\begin{ruledtabular}
\begin{tabular}{lp{0.70\columnwidth}}
(P0) & $L_w\le\sigma^{-1/2}$ (the window's Gaussian envelope in $w$ is not wider than $\sigma^{-1/2}$; no other condition on $L_w$)\\
(P1) & removed in v1.4: $f$ may have sine components, $f=1+\sum|\hat a_n|\cos(nt-\varphi_n)$ (App.~\ref{app:FN}; up to v1.3, $f$ even)\\
(P2) & $|\hat a_n|\le|\hat a_1|$: the fundamental is the largest harmonic\\
(P3) & $\Qb\le\tfrac14$ (blurred contrast)\\
(P4) & $\varrho=s^2/(2L^2)\le\tfrac18$ (window narrower than the envelope)\\
(P5) & $x=(k'\sigma)^2\ge1$ (period within the blur's reach); not needed for profiles with at most $20$ harmonics (any phases; App.~\ref{app:FN}) and for even strips (App.~\ref{app:P5})\\
(P6) & every state of the class has $|\varepsilon(\rho)|\le\bar\varepsilon$, $\bar\varepsilon\le0.1$ (paper~E's remainder, a state quantity); the comb $S'_C$ satisfies it for every $\bar\varepsilon\ge5.62\times10^{-7}$ (Theorem~\ref{thm:G1})\\
(P7) & $\sigma\le10^{-24}$ (both certificates cover $10^{-100000}\le\sigma\le10^{-24}$ continuously; below, the lower side by Lemma~\ref{lem:BH2} and the upper side by Theorem~\ref{thm:AZ1P})\\
\hline
\emph{not covered} & profiles whose fundamental is not the largest harmonic (a window of a shorter period);
$x<1$ for profiles with $21$ or more harmonics (separated lines in a window); windows with $L_w>\sigma^{-1/2}$\\
\end{tabular}
\end{ruledtabular}
\end{table}

\begin{theorem}[periodic windows: the logarithmic law from both sides] \caproofst{\code{bh\_cert\_cells.py}, \code{bf\_cert\_upper.py}, \code{bt\_eps.py}}
\label{thm:BF1}
Let $g_{\min}^P(\sigma)$ be the infimum of $g$ over the periodic windows that satisfy (P0), (P2)--(P7), together with the windows with at most $20$
harmonics that satisfy (P0), (P2)--(P7) except (P5). Then for every $\sigma$ with $10^{-100000}\le\sigma\le10^{-24}$
\begin{equation}
\begin{aligned}&(1-\bar\varepsilon)\Big(1-1.7\,\frac{\ln\ls}{\ls}\Big)\ \le\ g_{\min}^P\sqrt{\ls}\\ &\quad\le\ g(S'_C(\sigma))\sqrt{\ls}\\ &\quad\le\ \big(1+|\varepsilon(S'_C(\sigma))|\big)\Big(1+6.0\,\frac{\ln\ls}{\ls}\Big),\end{aligned}
\end{equation}
where $\bar\varepsilon\le0.1$ is the class bound of (P6), the left side is Theorem~\ref{thm:AY1P} ($421$ cells), the right side is Theorem~\ref{thm:AZ1P} on
the explicit comb $S'_C(\sigma)$ ($1464$ cells), and $|\varepsilon(S'_C(\sigma))|\le5.62\times10^{-7}$ (Theorem~\ref{thm:G1}). The middle inequality uses that $S'_C(\sigma)$ is in
the class; it holds for every $\bar\varepsilon\ge5.62\times10^{-7}$ without further condition (v1.4; up to v1.3 it was conditional on $|\varepsilon(S'_C)|\le\bar\varepsilon$).
\end{theorem}
\begin{proof}
The lower side is Theorem~\ref{thm:AY1P}, with Theorems~\ref{thm:E6} and~\ref{thm:F8} for the states with $x<1$. For the upper side, the comb sequence lies in the class on the whole range: (P0) since $L_w=8k_0^{-4}e^{-k_0^2}<1$; (P2) by construction;
$\Qb\le0.0056$, $\varrho\le3.5\times10^{-12}$ and $x\ge11.76$ (interval arithmetic on the same $1464$ cells, \code{bg\_members.py}). So (P3)--(P5) hold
(the former condition (P8) is no longer an assumption, Lemma~\ref{lem:BJ1}; the comb satisfies it anyway, $F_c\ge(1-|\varepsilon(S'_C)|)(1-1.7\times10^{-13})$, \code{bh\_members\_fc.py});
(P6) for the comb is Theorem~\ref{thm:G1}. Theorem~\ref{thm:AZ1P} then bounds $g$ on members of the class.
\end{proof}

\begin{corollary}[the limit] \label{cor:limit}
Let $\bar\varepsilon\in[5.62\times10^{-7},0.1]$ be the class bound of (P6). Then
\begin{equation*}
1-\bar\varepsilon\ \le\ \liminf_{\sigma\to0}g_{\min}^P\sqrt{\ls}\ \le\ \limsup_{\sigma\to0}g_{\min}^P\sqrt{\ls}\ \le\ 1 .
\end{equation*}
(i) \proofst{} The lower limit (Lemma~\ref{lem:BH2}). (ii) \caproofst{Theorems~\ref{thm:AZ1P} and~\ref{thm:G1}} The upper limit, and for the comb itself
$\lim_{\sigma\to0}g(S'_C(\sigma))\sqrt\ls=1$ without any condition. So the limit $1$ holds in the following sense, and in no stronger one: on every class with a
remainder bound $\bar\varepsilon$, the smallest gap times $\sqrt{\ln(1/\sigma)}$ ends between $1-\bar\varepsilon$ and $1$, and both ends tend to $1$ as
$\bar\varepsilon\to0$. What is not proved is a bound on paper~E's remainder $\varepsilon(\rho)$ over a class (the role of (P6)): the lower limit of the
infimum over all periodic windows, without a class bound, is open (on the states that approach the infimum $|\varepsilon|$ is numerically below $10^{-15}$).
\end{corollary}
\begin{proof}
(ii) For every periodic window $g\le(1+|\varepsilon|)U_P$; along $S'_C$, $U_P\sqrt\ls\to1$ (Theorem~\ref{thm:AZ1P}) and $\varepsilon(S'_C)\to0$ (Theorem~\ref{thm:G1}), so
$\limsup g(S'_C)\sqrt\ls\le1$; and Lemma~\ref{lem:BH2} applied to the classes $\{S'_C(\sigma'):\sigma'\le\sigma_0\}$, whose remainder bound $\sup_{\sigma'\le\sigma_0}|\varepsilon(S'_C(\sigma'))|$
tends to $0$ as $\sigma_0\to0$ (Theorem~\ref{thm:G1}), gives $\liminf g(S'_C)\sqrt\ls\ge1$. For every $\sigma\le10^{-24}$, $S'_C(\sigma)$ satisfies (P0), (P2)--(P5) and (P7) (\code{bg\_members.py} for $k_0\le240$, the asymptotic forms of
App.~\ref{app:cert} for $k_0\ge240$) and (P6) (Theorem~\ref{thm:G1}), so $g_{\min}^P\le g(S'_C)$.
\end{proof}

\paragraph*{Reading (observer).}
Whatever the shape of the teeth, the observer sees the window through its fundamental. The lower side says that the ripple cannot be finer
than $k_0^2\approx\tfrac14\ln(1/\sigma)$; the upper side says that the sharpest admissible comb reaches that limit. Sharpening the teeth
doubles the fundamental's contrast, which is worth dividing the blur by $256$ --- and changes $\ln(1/\sigma)$ only by $\ln256$, so the
constant stays $1$.
